\documentclass[14pt,a4paper]{extarticle}
% =============================================================
% preamble.tex — оформление отчётов по СТП БГУИР 01-2024
% =============================================================

% Шрифты: под pdflatex — Times через tempora + T2A,
% под xelatex/tectonic — системный Times New Roman (запасной вариант STIX Two Text)
\usepackage{iftex}
\ifPDFTeX
  \usepackage[utf8]{inputenc}
  \usepackage[T2A]{fontenc}
  \usepackage{tempora}
\else
  \usepackage{fontspec}
  \IfFontExistsTF{Times New Roman}
    {\setmainfont{Times New Roman}}
    {\setmainfont{STIX Two Text}}
  \IfFontExistsTF{DejaVu Sans Mono}
    {\setmonofont{DejaVu Sans Mono}[Scale=0.82]}
    {\setmonofont{DejaVuSansMono}[Extension=.ttf, BoldFont=*-Bold,
                                  ItalicFont=*-Oblique, Scale=0.82]}
\fi
\usepackage[russian]{babel}

% Поля: слева 30 мм, справа 15 мм, сверху и снизу 20 мм
\usepackage{geometry}
\geometry{left=30mm, right=15mm, top=20mm, bottom=20mm}

\usepackage{setspace}
\singlespacing

\pretolerance=1000
\tolerance=2000
\emergencystretch=3em

\usepackage{indentfirst}
\setlength{\parindent}{1.25cm}

% Номер страницы — снизу справа
\usepackage{fancyhdr}
\pagestyle{fancy}
\fancyhf{}
\fancyfoot[R]{\thepage}
\renewcommand{\headrulewidth}{0pt}
\fancypagestyle{plain}{%
  \fancyhf{}%
  \fancyfoot[R]{\thepage}%
  \renewcommand{\headrulewidth}{0pt}%
}

\usepackage{amsmath,amssymb}
\usepackage{graphicx}
\usepackage{float}
\usepackage{booktabs}
\usepackage{tabularx}
\usepackage{multirow}
\usepackage{longtable}
\usepackage{array}

\usepackage{caption}
\captionsetup[figure]{labelsep=endash, justification=centering,
                      font={normalsize}, name=Рисунок}
\captionsetup[table]{labelsep=endash, justification=raggedright,
                     singlelinecheck=false, font={normalsize}, name=Таблица}

\usepackage{titlesec}
\titleformat{\section}
  {\newpage\normalfont\fontsize{14}{18}\selectfont\bfseries}
  {\thesection}{0.5em}{\MakeUppercase}
\titlespacing*{\section}{1.25cm}{0pt}{14pt}
\titleformat{\subsection}
  {\normalfont\fontsize{14}{18}\selectfont\bfseries}
  {\thesubsection}{0.5em}{}
\titlespacing*{\subsection}{1.25cm}{14pt}{14pt}
\titleformat{\subsubsection}
  {\normalfont\fontsize{14}{18}\selectfont\bfseries}
  {\thesubsubsection}{0.5em}{}
\titlespacing*{\subsubsection}{1.25cm}{14pt}{14pt}

\newcommand{\sectioncenter}[1]{%
  \newpage
  \phantomsection
  \addcontentsline{toc}{section}{#1}
  \begin{center}\textbf{\MakeUppercase{#1}}\end{center}
  \vspace{14pt}
}

\usepackage{tocloft}
\renewcommand{\cftsecfont}{\normalfont}
\renewcommand{\cftsecpagefont}{\normalfont}
\renewcommand{\cftsubsecfont}{\normalfont}
\renewcommand{\cftsubsecpagefont}{\normalfont}
\renewcommand{\cftsecleader}{\cftdotfill{\cftdotsep}}
\renewcommand{\cftdotsep}{1}
% babel переопределяет заголовки после преамбулы, поэтому правим через \addto
\addto\captionsrussian{\renewcommand{\contentsname}{\hfill\textbf{СОДЕРЖАНИЕ}\hfill}}
\renewcommand{\cftaftertoctitle}{\hfill}

\usepackage{enumitem}
\setlist[itemize]{noitemsep, topsep=0pt, leftmargin=1.25cm, label=--}
\setlist[enumerate]{noitemsep, topsep=0pt, leftmargin=1.25cm}

% Листинги кода и вывода программы
\usepackage{listings}
\usepackage{fancyvrb}
\lstset{
  basicstyle=\ttfamily\footnotesize,
  frame=single,
  breaklines=true,
  showstringspaces=false,
  inputencoding=utf8,
  extendedchars=true,
  columns=fullflexible,
  keepspaces=true,
  captionpos=t
}
\ifPDFTeX
\lstset{
  literate={а}{{\cyra}}1 {б}{{\cyrb}}1 {в}{{\cyrv}}1 {г}{{\cyrg}}1
           {д}{{\cyrd}}1 {е}{{\cyre}}1 {ж}{{\cyrzh}}1 {з}{{\cyrz}}1
           {и}{{\cyri}}1 {й}{{\cyrishrt}}1 {к}{{\cyrk}}1 {л}{{\cyrl}}1
           {м}{{\cyrm}}1 {н}{{\cyrn}}1 {о}{{\cyro}}1 {п}{{\cyrp}}1
           {р}{{\cyrr}}1 {с}{{\cyrs}}1 {т}{{\cyrt}}1 {у}{{\cyru}}1
           {ф}{{\cyrf}}1 {х}{{\cyrh}}1 {ц}{{\cyrc}}1 {ч}{{\cyrch}}1
           {ш}{{\cyrsh}}1 {щ}{{\cyrshch}}1 {ъ}{{\cyrhrdsn}}1
           {ы}{{\cyrery}}1 {ь}{{\cyrsftsn}}1 {э}{{\cyrerev}}1
           {ю}{{\cyryu}}1 {я}{{\cyrya}}1 {ё}{{\cyryo}}1
           {А}{{\CYRA}}1 {Б}{{\CYRB}}1 {В}{{\CYRV}}1 {Г}{{\CYRG}}1
           {Д}{{\CYRD}}1 {Е}{{\CYRE}}1 {Ж}{{\CYRZH}}1 {З}{{\CYRZ}}1
           {И}{{\CYRI}}1 {Й}{{\CYRISHRT}}1 {К}{{\CYRK}}1 {Л}{{\CYRL}}1
           {М}{{\CYRM}}1 {Н}{{\CYRN}}1 {О}{{\CYRO}}1 {П}{{\CYRP}}1
           {Р}{{\CYRR}}1 {С}{{\CYRS}}1 {Т}{{\CYRT}}1 {У}{{\CYRU}}1
           {Ф}{{\CYRF}}1 {Х}{{\CYRH}}1 {Ц}{{\CYRC}}1 {Ч}{{\CYRCH}}1
           {Ш}{{\CYRSH}}1 {Щ}{{\CYRSHCH}}1 {Ъ}{{\CYRHRDSN}}1
           {Ы}{{\CYRERY}}1 {Ь}{{\CYRSFTSN}}1 {Э}{{\CYREREV}}1
           {Ю}{{\CYRYU}}1 {Я}{{\CYRYA}}1 {Ё}{{\CYRYO}}1
           {—}{{---}}1 {–}{{--}}1 {…}{{\ldots}}1
}
\fi

% Подписи листингов — в том же стиле, что таблицы: «Листинг X.Y — Название»
\captionsetup[lstlisting]{labelsep=endash, justification=raggedright,
                          singlelinecheck=false, font={normalsize}}
\renewcommand{\lstlistingname}{Листинг}
\AtBeginDocument{\numberwithin{lstlisting}{section}}

\lstdefinestyle{python}{language=Python, keywordstyle=\bfseries}
\lstdefinestyle{shell}{language=bash, keywordstyle=\mdseries}
\lstdefinestyle{plain}{language={}, basicstyle=\ttfamily\scriptsize}

% Абзац после листинга должен сохранять отступ 1,25 см (СТП 2.1.1)
\makeatletter
\AtBeginDocument{\let\@endpetrue\@endpefalse}
\makeatother

% Подпись к блоку verbatim: \captionof вне плавающего окружения
% глобально обнуляет \parindent, поэтому отступ восстанавливается вручную
\newcommand{\listingcaption}[2]{%
  \captionof{lstlisting}{#1}\label{#2}%
  \setlength{\parindent}{1.25cm}%
}

\usepackage{hyperref}
\hypersetup{colorlinks=true, linkcolor=black, urlcolor=blue, citecolor=black}

\numberwithin{figure}{section}
\numberwithin{table}{section}
\numberwithin{equation}{section}

% Титульный лист отчёта по лабораторной работе
\newcommand{\makelabtitle}[2]{%
  \begin{titlepage}
  \thispagestyle{empty}
  \begin{center}

  МИНИСТЕРСТВО ОБРАЗОВАНИЯ РЕСПУБЛИКИ БЕЛАРУСЬ

  \vspace{0.5em}
  Учреждение образования

  \vspace{0.5em}
  \textbf{<<Белорусский государственный университет информатики и радиоэлектроники>>}

  \vspace{2em}
  Факультет компьютерных систем и сетей

  \vspace{0.5em}
  Кафедра программного обеспечения информационных технологий
  \end{center}

  \vspace{3em}
  \begin{center}
  \textbf{ОТЧЕТ}

  \vspace{0.5em}
  по лабораторной работе №#1

  \vspace{0.5em}
  по дисциплине

  \vspace{0.5em}
  <<Информационные и технические средства защиты авторских прав>>

  \vspace{0.5em}
  <<#2>>
  \end{center}

  \vspace{5em}
  \begin{flushright}
  \begin{minipage}{0.5\textwidth}
  \textbf{Выполнил:}\\
  Лебедевич Артем Владимирович\\
  магистрант группы 556301

  \vspace{0.6em}
  \textbf{Преподаватель:}\\
  Ярмолик Вячеслав Николаевич\\
  доктор технических наук\\
  профессор кафедры ПОИТ
  \end{minipage}
  \end{flushright}

  \vfill
  \begin{center}Минск, 2026~г.\end{center}
  \end{titlepage}
}

% Титульный лист реферата
\newcommand{\makerefattitle}[1]{%
  \begin{titlepage}
  \thispagestyle{empty}
  \begin{center}

  МИНИСТЕРСТВО ОБРАЗОВАНИЯ РЕСПУБЛИКИ БЕЛАРУСЬ

  \vspace{0.5em}
  Учреждение образования

  \vspace{0.5em}
  \textbf{<<Белорусский государственный университет информатики и радиоэлектроники>>}

  \vspace{2em}
  Факультет компьютерных систем и сетей

  \vspace{0.5em}
  Кафедра программного обеспечения информационных технологий
  \end{center}

  \vspace{3em}
  \begin{center}
  \textbf{РЕФЕРАТ}

  \vspace{0.5em}
  по дисциплине

  \vspace{0.5em}
  <<Информационные и технические средства защиты авторских прав>>

  \vspace{1.5em}
  на тему

  \vspace{0.5em}
  \textbf{<<#1>>}
  \end{center}

  \vspace{5em}
  \begin{flushright}
  \begin{minipage}{0.5\textwidth}
  \textbf{Выполнил:}\\
  Лебедевич Артем Владимирович\\
  магистрант группы 556301

  \vspace{0.6em}
  \textbf{Преподаватель:}\\
  Ярмолик Вячеслав Николаевич\\
  доктор технических наук\\
  профессор кафедры ПОИТ
  \end{minipage}
  \end{flushright}

  \vfill
  \begin{center}Минск, 2026~г.\end{center}
  \end{titlepage}
}


\begin{document}

\makelabtitle{2}{Криптосистема с открытым ключом RSA}

\tableofcontents

\section{Цель работы}

Изучить принципы построения криптосистем с открытым ключом и алгоритм RSA.
В соответствии с вариантом №1 (нечётный номер в списке группы) реализовать шифратор
и дешифратор RSA, использующие алгоритм быстрого возведения в степень по модулю, а
также вычисление пары ключей по заданным значениям $p$, $q$ и $e$ с помощью
расширенного алгоритма Евклида. Проверить корректность реализации на примере из
методического пособия.

\section{Теоретические сведения}

\subsection{Криптосистемы с открытым ключом}

В симметричных шифрах, рассмотренных в лабораторной работе №1, отправитель и получатель
используют один и тот же секретный ключ, поэтому до начала переписки ключ необходимо
передать по защищённому каналу. Криптосистемы с открытым ключом (асимметричные)
устраняют этот недостаток: каждый участник генерирует пару ключей --- открытый $K_o$
и секретный $K_c$. Данные, зашифрованные открытым ключом, можно расшифровать только
секретным, причём вычислить $K_c$ по известному $K_o$ практически невозможно.
Следовательно, открытый ключ может распространяться по любым открытым каналам.

Основой таких систем служат односторонние функции с секретом: функции, которые легко
вычислить, трудно обратить, но легко обратить при знании дополнительной секретной
информации. Для RSA такой функцией является возведение в степень по модулю составного
числа $r = p \cdot q$, а секретом --- разложение $r$ на простые множители.

\subsection{Генерация ключей RSA}

Пара ключей RSA создаётся следующим образом:
\begin{enumerate}
\item Выбираются два секретных простых числа $p$ и $q$, близких по величине.
\item Вычисляется модуль $r = p \cdot q$.
\item Вычисляется функция Эйлера $\varphi(r) = (p - 1)(q - 1)$ --- количество чисел,
меньших $r$ и взаимно простых с ним.
\item Выбирается открытая экспонента $e$: $1 < e < \varphi(r)$, $\gcd(e, \varphi(r)) = 1$.
\item Вычисляется секретная экспонента $d$, удовлетворяющая условию
\begin{equation}
e \cdot d \bmod \varphi(r) = 1,
\end{equation}
то есть $d$ --- мультипликативное обратное к $e$ по модулю $\varphi(r)$.
\end{enumerate}

Открытым ключом является пара $K_o = (e, r)$, секретным --- $K_c = (d, r)$. Значения
$p$, $q$, $\varphi(r)$ и $d$ хранятся в секрете.

\subsection{Шифрование и дешифрование}

Сообщение $M$ разбивается на блоки $m_1, m_2, \dots$, каждый из которых представляет
собой число из диапазона $0 \dots r - 1$. Шифрование открытым ключом и дешифрование
секретным выполняются по формулам
\begin{equation}
c_i = m_i^{\,e} \bmod r, \qquad m_i = c_i^{\,d} \bmod r.
\end{equation}

Корректность следует из теоремы Эйлера. Поскольку $e d = 1 + k \varphi(r)$, для
$m$, взаимно простого с $r$, имеем
\begin{equation}
c^{d} = m^{ed} = m \cdot \left(m^{\varphi(r)}\right)^{k} \equiv m \cdot 1^k = m \pmod r.
\end{equation}
Для $m$, кратного $p$ или $q$, равенство доказывается отдельно по каждому модулю и
собирается по китайской теореме об остатках, поэтому дешифрование корректно для
любого $m < r$.

\subsection{Расширенный алгоритм Евклида}

Секретная экспонента находится расширенным алгоритмом Евклида (EUCLIDEX), который
одновременно с $\gcd(a, b)$ вычисляет коэффициенты $x$, $y$ соотношения Безу
\begin{equation}
x \cdot a + y \cdot b = \gcd(a, b).
\end{equation}

Если положить $a = \varphi(r)$, $b = e$ и числа взаимно просты, то
$y \cdot e \equiv 1 \pmod{\varphi(r)}$, то есть $y$ и есть искомое $d$. Отрицательное
$y$ приводится в диапазон прибавлением $\varphi(r)$. Алгоритм хранит две последние
тройки $(d, x, y)$ и на каждом шаге вычисляет новую по частному $q = d_0 \operatorname{div} d_1$.
Ход вычислений для примера из пособия ($\varphi(r) = 20$, $e = 7$) приведён в
таблице~\ref{tab:euclid}.

\begin{table}[H]
\caption{Расширенный алгоритм Евклида для $a = 20$, $b = 7$}
\label{tab:euclid}
\centering
\begin{tabular}{ccrrr}
\toprule
Шаг & $q$ & $d$ & $x$ & $y$ \\
\midrule
0 & --- & 20 & 1 & 0 \\
1 & --- & 7 & 0 & 1 \\
2 & 2 & 6 & 1 & $-2$ \\
3 & 1 & 1 & $-1$ & 3 \\
\bottomrule
\end{tabular}
\end{table}

Цикл останавливается при $d_1 = 1$: $(-1) \cdot 20 + 3 \cdot 7 = 1$, откуда $d = 3$.
Число итераций алгоритма Евклида не превышает $O(\log \varphi(r))$, поэтому вычисление
ключа выполняется мгновенно даже для 2048-битных модулей.

\subsection{Быстрое возведение в степень}

Непосредственное вычисление $m^e$ невозможно: для 1024-битных чисел результат
содержал бы порядка $2^{1024}$ цифр. Алгоритм быстрого возведения в степень
(fast\_exp) использует двоичное разложение показателя: пока показатель чётный, он
делится пополам, а основание возводится в квадрат; при нечётном показателе из него
вычитается единица и текущее основание домножается к результату. Все операции
выполняются по модулю $n$, поэтому промежуточные значения не превосходят $n^2$.
Количество умножений не превышает $2 \log_2 z$.

Ход вычисления $4^7 \bmod 33$ приведён в таблице~\ref{tab:exp}. Для показателя
$7 = 111_2$ выигрыш невелик, но для 1024-битного показателя требуется порядка
1500 умножений вместо $2^{1024}$.

\begin{table}[H]
\caption{Быстрое возведение в степень $4^7 \bmod 33$}
\label{tab:exp}
\centering
\begin{tabular}{rrrl}
\toprule
$z_1$ & $a_1$ & $x$ & Действие \\
\midrule
7 & 4 & 1 & начальные значения \\
6 & 4 & 4 & $z_1$ нечётно: $x = 1 \cdot 4$ \\
2 & 16 & 31 & $z_1 = 3$, $a_1 = 4^2$; затем $x = 4 \cdot 16 \bmod 33$ \\
0 & 25 & 16 & $z_1 = 1$, $a_1 = 16^2 \bmod 33$; затем $x = 31 \cdot 25 \bmod 33$ \\
\bottomrule
\end{tabular}
\end{table}

\subsection{Пример из методического пособия}

Для $p = 3$, $q = 11$ получаем $r = 33$, $\varphi(r) = 20$; при $e = 7$ расширенный
алгоритм Евклида даёт $d = 3$. Буквы английского алфавита кодируются числами
$A = 2, B = 3, \dots, Z = 27$ (значения 0 и 1 пропускаются, поскольку $0^e = 0$ и
$1^e = 1$ при любом ключе остались бы незашифрованными). Сообщение
$M = \text{<<CAT>>}$ принимает вид $\{4, 2, 21\}$:
\begin{equation*}
c_1 = 4^7 \bmod 33 = 16, \quad c_2 = 2^7 \bmod 33 = 29, \quad c_3 = 21^7 \bmod 33 = 21.
\end{equation*}
Дешифрование секретным ключом $K_c = (3, 33)$ восстанавливает исходные числа:
$16^3 \bmod 33 = 4$, $29^3 \bmod 33 = 2$, $21^3 \bmod 33 = 21$.

Заметим, что блок $m_3 = 21$ перешёл сам в себя. Такие неподвижные точки
($m^e \equiv m$) неизбежны в RSA; при малом модуле их доля заметна, при модуле
в тысячи бит вероятность попасть в неподвижную точку пренебрежимо мала.

\subsection{Криптостойкость RSA}

Стойкость RSA основана на двух трудных задачах. Восстановить $m$ по $c$ без ключа
означает извлечь корень степени $e$ по модулю составного числа. Вычислить $d$ по
открытому ключу означает найти $\varphi(r)$, что эквивалентно разложению $r$ на
множители. Ни для одной из этих задач не известно полиномиального алгоритма. Лучший
известный метод факторизации --- общий метод решета числового поля --- имеет
субэкспоненциальную сложность, и рекордная разложенная задача RSA-250 содержит
829 бит. В настоящее время рекомендуемая длина модуля составляет не менее 2048 бит.

\section{Реализация}

\subsection{Структура программы}

Программа написана на языке Python~3 без внешних зависимостей. Код второй
лабораторной работы находится в отдельном пакете \texttt{lab2}, общие
теоретико-числовые функции вынесены в \texttt{helpers/numtheory.py}, поскольку они
используются и в лабораторной работе №3 (листинг~\ref{lst:tree}).

\listingcaption{Структура программы}{lst:tree}
\begin{Verbatim}[fontsize=\footnotesize, frame=single, samepage=true]
main.py                CLI: команда rsa (keys, enc, dec), demo 2, gui 2
configs/cfg.py         p, q, e по умолчанию (пример из пособия)
helpers/numtheory.py   fast_exp, euclid_ex, inverse, тест Миллера-Рабина
lab2/rsa.py            ключи, шифрование, дешифрование, кодирование текста
lab2/runner.py         прогон примеров пособия и самопроверка
lab2/gui.py            окно: ключи, шифрование, пошаговый разбор
\end{Verbatim}

\subsection{Ключевые фрагменты кода}

Расширенный алгоритм Евклида (листинг~\ref{lst:euclid}) повторяет псевдокод
EUCLIDEX из пособия. Необязательный параметр \texttt{trace} собирает строки
таблицы~\ref{tab:euclid}, которые затем выводятся в демонстрации и в окне программы.

\begin{lstlisting}[style=python, caption={Расширенный алгоритм Евклида}, label={lst:euclid}]
def euclid_ex(a, b, trace=None):
    d0, d1 = a, b
    x0, x1 = 1, 0
    y0, y1 = 0, 1
    while d1 > 1:
        q, d2 = divmod(d0, d1)
        x2, y2 = x0 - q * x1, y0 - q * y1
        d0, d1, x0, x1, y0, y1 = d1, d2, x1, x2, y1, y2
        if trace is not None:
            trace.append(EuclidStep(q, d1, x1, y1))
    if d1 == 0:                      # b divides a: gcd is the previous remainder
        return x0, y0, d0
    return x1, y1, d1


def inverse(value, modulus, trace=None):
    _, y, g = euclid_ex(modulus, value % modulus, trace)
    if g != 1:
        raise ValueError(f"{value} is not invertible modulo {modulus}: gcd = {g}")
    return y % modulus               # a negative y1 is corrected by adding the modulus
\end{lstlisting}

Быстрое возведение в степень (листинг~\ref{lst:exp}) также реализовано строго по
пособию. Встроенная функция \texttt{pow(a, z, n)} языка Python не используется,
но служит эталоном при самопроверке: на 2000 случайных тройках с 30-значными числами
результаты совпали.

\begin{lstlisting}[style=python, caption={Быстрое возведение в степень по модулю}, label={lst:exp}]
def fast_exp(a, z, n, trace=None):
    a1, z1, x = a % n, z, 1 % n
    while z1 != 0:
        while z1 % 2 == 0:
            z1 //= 2
            a1 = a1 * a1 % n
        z1 -= 1
        x = x * a1 % n
        if trace is not None:
            trace.append(ExpStep(z1, a1, x))
    return x
\end{lstlisting}

Вычисление ключей (листинг~\ref{lst:keys}) проверяет все условия алгоритма: простоту
$p$ и $q$ (тестом Миллера--Рабина), их различие, диапазон $e$ и его взаимную простоту
с $\varphi(r)$. Нарушение любого условия приводит к понятному сообщению об ошибке,
а не к ключу, на котором дешифрование молча даст неверный результат.

\begin{lstlisting}[style=python, caption={Вычисление пары ключей}, label={lst:keys}]
def make_keys(p, q, e, trace=None):
    for name, value in (("p", p), ("q", q)):
        if not is_prime(value):
            raise ValueError(f"{name} = {value} is not prime")
    if p == q:
        raise ValueError("p and q must be different")
    phi = (p - 1) * (q - 1)
    if not 1 < e < phi:
        raise ValueError(f"e must satisfy 1 < e < phi(r) = {phi}")
    if gcd(e, phi) != 1:
        raise ValueError(f"e = {e} and phi(r) = {phi} are not coprime")
    return Keys(p, q, e, inverse(e, phi, trace))


def encrypt(blocks, e, r):
    _check_blocks(blocks, r)
    return [fast_exp(m, e, r) for m in blocks]


def decrypt(blocks, d, r):
    _check_blocks(blocks, r)
    return [fast_exp(c, d, r) for c in blocks]
\end{lstlisting}

Текст переводится в блоки одним из двух способов. Кодировка \texttt{letters}
соответствует пособию: $A = 2, \dots, Z = 27$, русские буквы (33, включая Ё)
продолжают таблицу: $\text{А} = 28, \dots, \text{Я} = 60$; пробелы и знаки препинания
отбрасываются. Кодировка \texttt{unicode} использует коды символов Unicode и позволяет
шифровать произвольный текст, но требует модуля $r > 1103$ для кириллицы. Если блок
не меньше $r$, программа сообщает об ошибке: такой блок нельзя восстановить, поскольку
дешифрование возвращает остаток по модулю $r$.

\subsection{Интерфейс командной строки}

Команда \texttt{rsa} работает в трёх режимах (листинг~\ref{lst:cli}). Ключи задаются
либо через $p$, $q$, $e$ (тогда $d$ вычисляется), либо напрямую парой $(e, r)$ для
шифрования и $(d, r)$ для дешифрования --- так, как их получает собеседник.

\begin{lstlisting}[style=shell, caption={Примеры запуска}, label={lst:cli}]
python3 main.py rsa keys -p 3 -q 11 -e 7     # r, phi(r), ход алгоритма Евклида, d
python3 main.py rsa enc -t CAT                # 16 29 21 (ключи из пособия)
python3 main.py rsa dec -d 3 -r 33 -t "16 29 21"            # CAT
python3 main.py rsa enc -p 61 -q 53 -e 17 -t "Привет"       # 678 1086 1350 ...
echo "Hi, мир!" | python3 main.py rsa enc --codec unicode -p 1009 -q 1013 -e 17
python3 main.py demo 2                        # все примеры с проверкой
python3 main.py gui 2                         # графическое окно
\end{lstlisting}

\subsection{Графическое приложение}

Окно лабораторной работы №2 (\texttt{python3 main.py gui 2}) содержит поля $p$, $q$,
$e$, кнопку вычисления ключей, поля открытого текста и шифротекста, выбор кодировки
и панель <<Как это работает>>. При вычислении ключей на панели выводятся $r$,
$\varphi(r)$ и таблица расширенного алгоритма Евклида, при шифровании или
дешифровании --- блоки до и после преобразования и ход быстрого возведения в степень
для первого блока. Интерфейс переключается между английским и русским языками.

Общая часть окна (поля ключей, панель разбора, переключатель языка) оформлена базовым
классом \texttt{KeyWindow}, который переиспользуется окном лабораторной работы №3.

\section{Результаты работы программы}

Команда \texttt{python3 main.py demo 2} выполняет пример из пособия, шифрование
русского текста и произвольного текста с 256-битным модулем, проверку отклонения
неверных параметров и самопроверку. Сводные результаты приведены в
таблице~\ref{tab:results}, полный вывод --- в листинге~\ref{lst:demo}.

\begin{table}[H]
\caption{Результаты работы программы}
\label{tab:results}
\centering
\begin{tabularx}{\textwidth}{Xll}
\toprule
Проверка & Результат & Эталон \\
\midrule
Секретная экспонента для $p = 3$, $q = 11$, $e = 7$ & $d = 3$ & совпадает \\
Шифрование <<CAT>> & 16 29 21 & совпадает \\
Дешифрование 16 29 21 & CAT & совпадает \\
Русский текст, $r = 3233$, $d = 2753$ & восстановлен & --- \\
Unicode-текст, 256-битный модуль & восстановлен & --- \\
\texttt{fast\_exp} против \texttt{pow}, 2000 троек & совпадение & --- \\
50 случайных пар ключей по 20 блоков & восстановлены & --- \\
\bottomrule
\end{tabularx}
\end{table}

Неверные параметры отклоняются с пояснением: $e = 5$ не взаимно просто с
$\varphi(33) = 20$, число 4 не является простым, $e = 25$ выходит за границу
$\varphi(r)$, а русские буквы при $r = 33$ не помещаются в диапазон блоков.

\listingcaption{Вывод программы}{lst:demo}
\VerbatimInput[fontsize=\scriptsize, frame=single]{listings/demo_output.txt}

\section{Анализ результатов}

Пример с русским текстом наглядно показывает главную слабость посимвольного
применения RSA: одинаковые буквы дают одинаковые блоки шифротекста. Буква <<А>>
сообщения <<Защита авторских прав>> четыре раза превратилась в число 1075, буква
<<В>> --- дважды в 1726. Фактически получен шифр простой замены с алфавитом из
чисел, который вскрывается частотным анализом так же, как шифр Цезаря из лабораторной
работы №1, несмотря на стойкость самого RSA. Кроме того, при открытом ключе
противник может просто зашифровать все 59 возможных букв и составить словарь.

Поэтому на практике RSA применяют иначе:
\begin{itemize}
\item несколько символов объединяются в один блок длиной почти $\log_2 r$ бит;
\item к блоку добавляется случайное дополнение (схема OAEP), так что одно и то же
сообщение каждый раз шифруется по-разному;
\item RSA шифрует не само сообщение, а случайный сеансовый ключ симметричного шифра
(AES), которым шифруются данные: возведение в степень с 2048-битным модулем в тысячи
раз медленнее симметричного шифрования.
\end{itemize}

Ключи из пособия ($r = 33$, $r = 3233$) пригодны лишь для иллюстрации: модуль такой
величины раскладывается перебором делителей до $\sqrt{r}$ за доли миллисекунды, после
чего $d$ вычисляется тем же расширенным алгоритмом Евклида. Даже 256-битный модуль из
демонстрации разлагается на персональном компьютере за минуты.

\sectioncenter{Заключение}

В ходе лабораторной работы изучены принципы криптографии с открытым ключом и
реализована криптосистема RSA по варианту №1: вычисление пары ключей по заданным
$p$, $q$, $e$ расширенным алгоритмом Евклида, шифратор и дешифратор на основе
быстрого возведения в степень по модулю. Программа предоставляет интерфейс командной
строки и графическое окно с пошаговым разбором вычислений. Корректность подтверждена
совпадением с примером пособия ($d = 3$, шифротекст 16 29 21), сравнением
\texttt{fast\_exp} со встроенной функцией \texttt{pow} и проверкой обратимости на
случайных ключах с модулем до 190 бит.

Практический вывод работы состоит в том, что стойкость алгоритма не переносится
автоматически на способ его применения: посимвольное шифрование превращает RSA в шифр
простой замены. Надёжная схема требует блоков максимальной длины, случайного
дополнения и модуля не менее 2048 бит, а для больших объёмов данных --- гибридной
схемы, где RSA передаёт только ключ симметричного шифра.

\sectioncenter{Список использованных источников}

\begin{enumerate}
\item Ярмолик, С.~В. Криптография и охрана коммерческой информации: метод. пособие по
выполнению лаб. работ / С.~В.~Ярмолик, В.~Н.~Ярмолик. --- Минск: БГУИР, 2010. --- 32~с.
\item Rivest, R. A method for obtaining digital signatures and public-key cryptosystems /
R.~Rivest, A.~Shamir, L.~Adleman // Communications of the ACM. --- 1978. --- Vol.~21,
No.~2. --- P.~120--126.
\item Шнайер, Б. Прикладная криптография. Протоколы, алгоритмы, исходные тексты на
языке Си / Б.~Шнайер. --- М.: Триумф, 2002. --- 816~с.
\item Столлингс, В. Криптография и защита сетей: принципы и практика /
В.~Столлингс. --- М.: Вильямс, 2001. --- 672~с.
\end{enumerate}

\end{document}
